\section{Operational controls for automated trading}

\begin{qfmethod}[Operational controls for automated trading: reliability and recovery, unknown order state, audit trails, and the U.S. market-access rule]\label{meth:automated-trading-operational-controls}
\qffield{Statement} Harris describes the operational demands on automated trading systems of markets, ``as well as brokers and dealers'', as of 2002 \citep[chap.~``Floor Versus Automated Trading Systems'', sec.~27.7.1, pp.~549--551]{alpha-harris}. The systems must validate users, accept and process orders, report trades and report order status, reliably, securely and fast:
\begin{itemize}
\item \emph{reliability}: trading stops when the computer or network is down, and traders will not use unreliable systems. Builders therefore eliminate single points of failure and, since failures are inevitable, invest in recovering from interruptions: fault-tolerant hardware; redundant computers and network connections through different vendors and physical paths; backup power; a remote ``hot disaster recovery site'' to which activity can switch at once; significant excess capacity for surges in demand; quality controls in software development; redundant data backups; and thorough testing by independent quality-assurance groups (pp.~549--550);
\item \emph{security}: encryption and user authentication, so that hackers cannot tamper with systems and impersonators cannot trade in others' names (p.~550);
\item \emph{speed}, because traders without the latest information decide badly (pp.~550--551).
\end{itemize}
An \emph{audit trail} records the submission and disposition of every order, with detailed information on everything that happens to it. Regulators use it to decide whether traders have violated trading rules; floor rules require time stamps when orders are received and filled, sequential recording and immediate reporting; automated systems produce complete audit trails of all activity within them (sec.~27.1.1, pp.~544--545).

When a connection fails, Harris's examples show three distinct risks (box, p.~550). A trader whose accepted limit order may or may not have filled cannot safely trade elsewhere: if she does so without knowing, she risks trading twice. An arbitrageur who has bought one leg cannot sell the other, and is left unhedged. A speculator cannot exit a losing position; a stop order placed immediately after the purchase would have represented him while he could not communicate.

In the United States, Rule 15c3-5 requires a broker or dealer with market access, or that provides others with access to an exchange or alternative trading system, to establish, document and maintain risk-management controls and supervisory procedures \citep[17 CFR 240.15c3-5(b)]{cfr-240-15c3-5}. These are reasonably designed, among other things:
\begin{itemize}
\item to prevent orders that exceed pre-set credit or capital thresholds ((c)(1)(i));
\item to prevent erroneous orders exceeding price or size parameters, singly or over a short period, or ``that indicate duplicative orders'' ((c)(1)(ii));
\item to restrict access to trading systems and technology to pre-approved, authorised persons and accounts ((c)(2)(iii));
\item to assure that surveillance personnel receive immediate post-trade execution reports ((c)(2)(iv)).
\end{itemize}
The controls must be under the firm's direct and exclusive control, with a limited contractual allocation of some regulatory controls to broker-dealer customers ((d)); they must be reviewed regularly, at least annually, with an annual certification by the chief executive officer ((e)).

\qffield{Conditions} Interpretation, a synthesis of the passages above and not a procedure stated in them: after a failure the states to separate are these.
\begin{itemize}
\item \emph{Unknown order state}: an order may be unacknowledged, live, partly filled, filled, cancelled or rejected. Venue messages report these transitions (acknowledgements, rejections, cancellations and fills; \hyperref[def:automated-trading-architecture]{automated trading architecture}), and the trader's view is stale until they arrive.
\item \emph{Unknown fill state and position}: what matters for risk is the quantity filled and the resulting position, including the unhedged leg of a multi-leg trade.
\item \emph{Reconciliation}: establishing state from authoritative records, such as venue or broker execution reports and post-trade reports, before acting.
\item \emph{Retry}: resubmitting before reconciliation is how a trader trades twice. For U.S. broker-dealers with market access, rejecting orders that indicate duplicative orders is a required pre-trade control, but such a filter is a backstop, not a reconciliation.
\item \emph{Cancellation}: a cancel sent while state is unknown may cross with a fill, so its outcome also needs confirmation.
\item \emph{Recovery}: redundancy, backups and disaster-recovery sites restore service; they do not by themselves establish what happened to orders in flight.
\end{itemize}

\qffield{Practice} Interpretation: the advice to prefer safe state reconciliation over blind resubmission follows from Harris's ``trading twice'' and is consistent with the duplicative-order control of Rule 15c3-5, which states the control, not the reconciliation procedure. Harris's stop-order example is a way to leave a resting instruction that protects a position while disconnected; it is not a substitute for knowing the order's state. The Knight Capital order (\nolinkurl{wide-knight-capital}) is summarised in \hyperref[ev:quant-implementation-failures]{implementation failures}.

\qffield{Pitfalls} Keep three kinds of statement apart.
\begin{itemize}
\item \emph{Description}: Harris's lists are descriptions of practice in 2002 markets, not requirements.
\item \emph{Law}: Rule 15c3-5 is U.S. law for broker-dealers with market access, as in its current eCFR text. It does not mention kill switches, reconciliation or ``drop copies''; the immediate post-trade execution reports it requires are the nearest provision.
\item \emph{EU}: the algorithmic-trading obligations are in \hyperref[rule:mifid-rts6-algorithmic-controls]{RTS~6} (\nolinkurl{wide-mifid-rts6}). U.K. supervisory statements, industry guidance (\nolinkurl{w2-fia-automated-controls}), Regulation SCI and audit-trail rules such as the U.S. consolidated audit trail are not covered here.
\end{itemize}
The cited passages support less than common practice:
\begin{itemize}
\item reliability: Harris names validating users; drop copies, replayable feeds, tamper-resistant logs and order validation are additions;
\item timing: Harris states a need for speed, not bounded delays for each stage;
\item market quality: Harris says only that traders will not use unreliable systems and that trading stops when systems are down;
\item audit trails: Harris names front running, inappropriate order exposure, fraudulent trade assignment and prearranged trading as the floor scandals that audit trails help deter; allocation, linking private messages to public market state, and tests of best execution, front running, allocation and unauthorised trading are additions.
\end{itemize}

\qffield{See also} \hyperref[def:automated-trading-architecture]{trading architecture} (venue and recorded messages, model checks); \hyperref[def:floor-automated-trading-systems]{floor against automated systems}; \hyperref[def:best-execution]{best execution}.

\qffield{Further reading} The adopting release of Rule 15c3-5 is \nolinkurl{wide-sec-15c3-5}.

\end{qfmethod}

\begin{qfrule}[EU organisational requirements for algorithmic trading (RTS~6): testing and deployment, annual self-assessment, kill functionality, pre- and post-trade controls, real-time monitoring and business continuity]\label{rule:mifid-rts6-algorithmic-controls}
\qffield{Statement} Commission Delegated Regulation (EU) 2017/589 of 19 July 2016, supplementing Article~17 of MiFID~II (Directive 2014/65/EU), sets the organisational requirements of investment firms engaged in algorithmic trading; it applies from 3 January 2018 (Art.~29) \citep[OJ L 87, 31.3.2017, pp.~417ff.]{wide-mifid-rts6}. Its main obligations:
\begin{itemize}
\item \emph{Development and deployment}: clearly delineated methodologies to develop and test any algorithmic trading system, algorithm or strategy before deployment or substantial update, authorised by a person designated by senior management, ensuring among other things that it ``does not behave in an unintended manner'' (Art.~5); conformance testing with venues (Art.~6), separate testing environments (Art.~7), and predefined limits before deployment on the number of instruments, the price, value and number of orders, strategy positions and the number of venues (Art.~8).
\item \emph{Post-deployment}: an annual self-assessment and validation report drawn up by the risk-management function, audited by internal audit where it exists and approved by senior management, with deficiencies remedied (Art.~9); within it, stress tests of messaging and trade volumes at twice the highest of the previous six months, run without affecting production (Art.~10); review of material changes by a designated person (Art.~11).
\item \emph{Kill functionality}: the ability ``to cancel immediately, as an emergency measure, any or all of its unexecuted orders submitted to any or all trading venues'', and to identify which algorithm and which trader, desk or client is responsible for each order (Art.~12).
\item \emph{Business continuity}: documented arrangements covering adverse scenarios, back-up sites where appropriate, use of the kill functionality, shutdown of algorithms ``without creating disorderly trading conditions'' and management of outstanding orders and positions, reviewed and tested annually (Art.~14).
\item \emph{Pre-trade controls} for all instruments: price collars, maximum order values and volumes, and maximum message limits, with all orders sent included immediately in the limits; repeated-execution throttles after which the system is disabled until re-enabled by designated staff; market and credit risk limits; blocking of unauthorised traders and of orders that compromise risk thresholds; and an exceptional, temporary override procedure verified by risk management (Art.~15).
\item \emph{Real-time monitoring} of all algorithmic activity under the firm's trading code by the trader in charge and by an independent risk function, with alerts ``generated within five seconds after the relevant event'' and reachable staff (Art.~16).
\item \emph{Post-trade controls}, operated continuously, including reconciliation of the firm's ``own electronic trading logs with information about its outstanding orders and risk exposures as provided by the trading venues'', brokers, clearing members and data providers, ``in real-time where the aforementioned market participants provide the information in real-time'', and real-time calculation of outstanding exposure (Art.~17(1)--(3)).
\end{itemize}

\qffield{Conditions} Jurisdiction and scope: EU investment firms engaged in algorithmic trading under MiFID~II Art.~17; Articles 19--23 cover direct electronic access providers and Articles 24--26 general clearing members. As of October 2026 there is no amending act and one corrigendum (28 July 2017, German language version only), so the English text is the text in force. Whether Directive 2024/790 changed Article~17 itself, or ESMA has proposed revised standards, is unverified.

\qffield{Practice} Interpretation: the regulation states \emph{what} controls must exist and who signs them off, not their parameters; collars, limits and throttles are calibrated by the firm. Compared with the U.S. market-access rule (\hyperref[meth:automated-trading-operational-controls]{operational controls}), RTS~6 names kill functionality, real-time reconciliation with venues and brokers, and annual stress tests explicitly; Rule 15c3-5 names none of these.

\qffield{Pitfalls} The text is the Official Journal version; Annexes~I (self-assessment criteria) and~II are not covered. For national competent authorities' and ESMA's guidance and supervisory reviews see, for the U.K., \nolinkurl{wide-fca-algo-review}. The U.K. retains its own version of this regulation after 2020.

\qffield{See also} \hyperref[meth:automated-trading-operational-controls]{operational controls for automated trading}; \hyperref[mech:order-state-reconciliation]{order-state reconciliation}; \hyperref[ev:quant-implementation-failures]{implementation failures}; \hyperref[mech:circuit-breakers-trading-halts]{circuit breakers}.

\qffield{Further reading} \nolinkurl{w2-fia-automated-controls} (industry guidance).

\end{qfrule}

\begin{qfmechanism}[Reconciling order state after a failure: FIX status requests, execution reports, resend and possible-duplicate flags, and the EU reconciliation duty]\label{mech:order-state-reconciliation}
\qffield{Statement} The FIX application layer gives a client three tools for establishing what happened to its orders \citep[FIX Latest, EP312, Orchestra repository of 8 September 2026]{s6-fix}:
\begin{itemize}
\item \emph{OrderStatusRequest} (35=H) ``is used by the order submitter to generate an ExecutionReport(35=8) message with ExecType(150)=I (Order Status) back from the recipient''; OrdStatusReqID (790) can identify the request;
\item \emph{OrderMassStatusRequest} (35=AF) requests the status of all orders matching criteria in MassStatusReqType (585), and an ExecutionReport with ExecType=I is returned ``for each order matching the criteria provided on the request with the order's status'';
\item the \emph{ExecutionReport} (35=8) itself, which acknowledges, confirms changes, relays status and fills and rejects orders, and carries ``the current state of the order as understood by the recipient (OrdStatus(39)) and the purpose of the message (ExecType(150))''. Its OrdStatus precedence table ranks Pending Cancel (11) and Pending Replace (10) above Done for Day, Filled, Canceled, Partially Filled and New; Pending Cancel ``DOES NOT INDICATE THAT THE ORDER HAS BEEN CANCELED''; an order ``can pass from New to Rejected status''; and Pending New is reported ``only \dots\ in response to an OrderStatusRequest(35=H)''.
\end{itemize}
The session layer recovers lost messages: \emph{ResendRequest} (35=2) initiates retransmission; resent messages carry PossDupFlag (43), which ``indicates possible retransmission of the message with the sequence number in the MsgSeqNum(34) field'', with OrigSendingTime (122); PossResend (97) marks a message that ``may have been sent earlier under another sequence number''; \emph{SequenceReset} (35=4) with GapFillFlag (123)\,$=$\,Y fills gaps over messages that will not be resent \citep[FIXT session layer, Orchestra file FIX.5.0SP2\_EP247 of 16 May 2019]{s6-fix}. In the EU, the firm ``shall reconcile its own electronic trading logs with information about its outstanding orders and risk exposures'' from venues, brokers, clearing members and data providers, in real time where they provide real-time information (RTS~6, Art.~17(3); \hyperref[rule:mifid-rts6-algorithmic-controls]{RTS~6}).

\qffield{Notation} Tag numbers in parentheses; 35 is MsgType. ClOrdID (11) is the client's order identifier and OrderID (37) the recipient's.

\qffield{Conditions} The documented method, assembled from those definitions (the arrangement is the entry's, each step is the source's): (1) restore the session and let the counterparty resend, treating PossDupFlag messages as possible duplicates of sequence numbers already processed and PossResend messages as possible duplicates of business content under new sequence numbers, so that a fill is applied once, keyed by its execution identifier; (2) request status, order by order (35=H) or in bulk (35=AF), and take OrdStatus and the cumulative and leaves quantities in the ExecType=I replies as the recipient's view; (3) compare with the internal order log and position, and with clearing or drop-copy records where available (the RTS~6 duty); (4) treat Pending Cancel and Pending Replace as unresolved, not as cancelled or replaced; (5) only then decide whether anything may be sent.

Step~(5) needs facts that reconciliation alone does not supply. A reconciled, terminal order is not thereby a candidate for resubmission: a fully filled order is terminal and reconciled and has nothing left to send. A replacement is justified only when three conditions hold together:
\begin{itemize}
\item the parent instruction (the intent) is still valid, not withdrawn by the strategy or a risk control;
\item every child order of that intent is reconciled and resolved;
\item authorised quantity remains, after deducting filled and still-fillable quantity (live leaves, plus the open quantity of unresolved or broken orders).
\end{itemize}
\emph{Complete for today is not ended.} ``Terminal'' here means that the order's lifetime has ended: Filled, Canceled, Rejected or Expired. Done for Day (OrdStatus~3) is ``Order not, or partially, filled; no further executions forthcoming for the trading day''. Expired is an order ``canceled in broker's system due to TimeInForce(59) instructions''. A GoodTillCancel order ``remains in effect until it is either executed or canceled''. LeavesQty ``could be 0'' when OrdStatus is Done for Day \citep[FIX Latest, EP312, OrdStatus precedence table, codes 59 and field 151]{s6-fix}.

FIX 4.4's state-change matrices, as restated in OnixS's mirror of Appendix~D, show the two cases. A day order ends Done for Day with LeavesQty~0 (A.1.b, ``Assuming day order''). A GTC order reports Done for Day with 8{,}000 of 10{,}000 still leaves, and is restated the next morning (ExecType~D, ExecRestatementReason ``GTC renewal/restatement'') before filling again (H.1.a) \citep[app.~D, matrices A.1.b and H.1.a]{onixs-fix44-appd}.

So Done for Day ends an order's lifetime only for a day order, and only where the venue's rules say no Expired or Canceled report follows. A good-till order (GTC, GTD, or TimeInForce unknown in the log) is dormant: it can fill in the next session for OrderQty less CumQty, whatever LeavesQty it reported. Its remaining quantity is released for replacement only after a confirmed cancellation or expiry. The replacement's size is that remaining quantity. This separation of order state, remaining authorised quantity and intent is the entry's own arrangement; FIX defines the order state, not the intent. Field-level rules for deduplication (ExecID, ExecRefID, corrections) are not covered here.

The state-change matrices are cited for Done for Day (A.1.b) and GTC renewal (H.1.a) from OnixS's FIX~4.4 restatement; the Orchestra files carry message and field definitions and the OrdStatus precedence table, not the matrices.

\qffield{Practice} The Knight failure of 1 August 2012 was in part a state failure of this kind: ``one part of Knight's order handling system recognized that the parent orders had been filled, this information was not communicated'' to the router, which kept sending child orders, and Knight had no control ``to compare orders leaving SMARS with those that entered it'' \citep[paras.~16 and~21, pp.~6--7]{wide-knight-capital} (\hyperref[ev:quant-implementation-failures]{implementation failures}). It covers a duplicate (PossDup) report applied once, a fully filled order that is \emph{not} resubmitted, a partially filled cancelled order (180 of 300 proposed when its intent is valid, nothing when it is withdrawn), an order in Pending Cancel, a missing execution found from a status reply, and a GTC order Done for Day with zero reported leaves.

That order is dormant, not terminal: 8{,}000 could still fill, nothing is proposed, it is renewed and partly filled the next day, and 7{,}000 is released only after a confirmed cancel. A day order Done for Day is ended only under a stated venue rule. It illustrates the logic and is not a test against a venue.

\qffield{Pitfalls} Venue behaviour varies: whether a venue supports 35=H or 35=AF, honours resend requests or answers only through drop copies is set in its rules of engagement, not by FIX. Interpretation: an ExecType=I reply is the recipient's view at the time it is sent; a fill in flight can still arrive after it.

\qffield{See also} \hyperref[meth:automated-trading-operational-controls]{operational controls}, whose unknown-state list this entry sources; \hyperref[mech:order-lifecycle]{order lifecycle}; \hyperref[rule:mifid-rts6-algorithmic-controls]{RTS~6}.

\qffield{Further reading} The Orchestra repository \nolinkurl{github.com/FIXTradingCommunity/orchestrations}, commit \texttt{cd24169}; the FIX specification volumes and the FIX Session Layer specification.

\end{qfmechanism}

\begin{qfevidence}[Implementation failures in quantitative trading: Knight Capital (1 August 2012) and AXA Rosenberg (2007--2010) in the regulators' orders, against a practitioner's account]\label{ev:quant-implementation-failures}
\qffield{Statement} \emph{Knight Capital.} On 1 August 2012 Knight's equity router SMARS, ``while processing 212 small retail orders'', routed millions of orders over 45 minutes and obtained ``over 4 million executions in 154 stocks for more than 397 million shares'', leaving net long positions of about \$3.5 billion in 80 stocks and net short positions of about \$3.15 billion in 74; Knight ``lost over \$460 million'' \citep[para.~1, p.~2]{wide-knight-capital}. New code for the NYSE Retail Liquidity Program repurposed a flag that had activated ``Power Peg'', a function unused since 2003 whose cumulative-quantity check had been moved in 2005 without retesting; one of eight servers did not receive the new code, without a second technician's review or written procedures; on that server the flag triggered Power Peg, which sent child orders ``without regard to the number of share executions'' (paras.~13--16, pp.~5--6).

Ninety-seven automated e-mails referring to ``Power Peg disabled'' before the open were not designed as alerts and not acted on (para.~19, pp.~6--7); there was no control comparing orders leaving SMARS with those entering it, no firm-wide capital threshold linked to automated order blocking, a post-trade monitor (PMON) relying on human monitoring without automated alerts, and no incident-response procedure; uninstalling the new code from the seven correct servers made matters worse (paras.~21--27, pp.~7--8). The SEC found violations of Rule 15c3-5, censured Knight and imposed a \$12 million penalty (paras.~9--10, pp.~3--4; sec.~IV, p.~18).

\emph{AXA Rosenberg.} An error introduced in April 2007 when linking a new risk model to the optimiser failed to scale inputs reported partly in decimals and partly in percentages, so the optimiser ``did not give the intended weight to common factor risks''; found in June 2009, it was concealed on a senior official's instruction, fixed in September--November 2009, disclosed to the SEC's examiners in late March 2010 and to clients on 15 April 2010; it affected 608 of 1{,}421 portfolios and caused \$216{,}806{,}864 of losses, which the respondents undertook to pay, with a \$25 million penalty \citep[paras.~1--3, 10--11, 33; sec.~IV, pp.~2--4, 10 and~16]{sec-axa-rosenberg-2011}.

\qffield{Conditions} Settled orders: findings are made ``without admitting or denying'' them and bind only the parties. Two cases, chosen because Narang cites them; they are case studies of controls, not estimates of how often such failures occur.

\qffield{Practice} Narang's account \citep[sec.~10.1.3, pp.~195--196]{s7-narang} differs from the orders in detail: he gives Knight's loss as ``more than \$400 million in a mere 30 minutes'' and attributes it to a bug that ``caused a dormant piece of software to come back online, multiplying order sizes''; the order gives over \$460 million in about 45 minutes and describes repeated child orders sent without regard to fills, not larger orders. His AXA figures (\$217 million, \$25 million) match the order; his ``about three years later, in April 2010'' measures from the introduction of the error (April 2007), not from its discovery (June 2009). His claims about Knight's later recapitalisation, a March 2011 ETF ``process error'' and AXA's assets under management are not in the orders .

For control design the orders point to: deployment review and removal of dead code, an output control on the router (orders out against orders in), automated linkage of exposure limits to order entry, alerts designed as alerts, and an incident procedure that does not change production blindly; RTS~6 names several of these as EU obligations (\hyperref[rule:mifid-rts6-algorithmic-controls]{RTS~6}, Arts.~8, 12, 15--17).

\qffield{Pitfalls} Reading the AXA case as a mere ``coding error'' omits the violation the order identifies: concealment and misleading statements to clients.

\qffield{See also} \hyperref[def:quant-strategy-risk-sources]{sources of strategy risk} (Narang's taxonomy); \hyperref[meth:automated-trading-operational-controls]{operational controls}; \hyperref[mech:order-state-reconciliation]{order-state reconciliation}.

\qffield{Further reading} The SEC's other technology-governance orders cited in the Knight order's n.~1 (NASDAQ 2013, NYSE 2012, EDGX 2011).

\end{qfevidence}
